We use the following masked generalization of X3C.
For the general masked system below, \(M\) denotes the number of equations.
In the X3C specialization, \(M=m\).
Let
\(M,n\geq1\), and consider the \emph{masked Boolean system}
\begin{equation}
 \mathbf A\boldsymbol\xi=\mathbf b,
 \qquad
 \boldsymbol\xi\in\{0,1\}^n,
 \qquad
 \boldsymbol\xi\leq\boldsymbol\mu,
 \label{eq:masked-system}
\end{equation}
where \(\mathbf A\in\{0,1\}^{M\times n}\),
\(\mathbf b\in\{0,1\}^M\), and
\(\boldsymbol\mu\in\{0,1\}^n\).

In this paper, we call \(\boldsymbol\mu\) the \emph{mask}.  The
constraint \(\boldsymbol\xi\leq\boldsymbol\mu\) permits \(\xi_k=1\) only when
\(\mu_k=1\).  The feasibility problem for masked Boolean systems is
\(\NP\)-complete,
since X3C is the special case obtained by setting
\[
 (\mathbf A,\mathbf b,\boldsymbol\mu)
 \mathrel{:=} (\mathbf H,\mathbf 1_m,\mathbf 1_n).
\]

We reduce this feasibility problem to exact decision-CVP on
principal ideal lattices.  Roughly speaking, for each Boolean vector
\(\boldsymbol\xi\leq\boldsymbol\mu\), we form a
sparse polynomial \(a_{\boldsymbol\xi}\) and construct a principal ideal
\((g)\), a target \(T_{\boldsymbol\mu}\), and a threshold
\(\Delta_{\boldsymbol\mu}\) with
\[
 \|g(y) a_{\boldsymbol\xi}(y^2)
       -T_{\boldsymbol\mu}(y)\|_2^2
 =\Delta_{\boldsymbol\mu}
  +4\|\mathbf A\boldsymbol\xi-\mathbf b\|_2^2.
\]
We then add coefficients at odd powers of \(y\).  If \(q\)
is not of the form \(a_{\boldsymbol\xi}(y^2)\) for any
\(\boldsymbol\xi\in\{0,1\}^n\) satisfying
\(\boldsymbol\xi\leq\boldsymbol\mu\), the resulting ideal element \(gq\)
has squared distance from \(T_{\boldsymbol\mu}\) greater than
\(\Delta_{\boldsymbol\mu}+4\).  Thus
the closest squared distance is \(\Delta_{\boldsymbol\mu}\) for a YES
instance and at least \(\Delta_{\boldsymbol\mu}+4\) for a NO instance.

\subsection{Encoding the Boolean equations}

We first choose coefficient positions \(\alpha_1,\ldots,\alpha_n\) whose
nonzero ordered differences are pairwise distinct.
Such positions are obtained from the classical
Erd{\H{o}}s--Tur{\'a}n Sidon-sequence
construction~\cite{ErdosTuran41}.

\begin{lemma}
\label{lem:sidon}
There exists a deterministic polynomial-time algorithm that, on input
\(n\geq1\), outputs integers
\[
 0=\alpha_1<\cdots<\alpha_n,
 \qquad
 \alpha_n=\Theta(n^2)\quad(n\geq2),
\]
such that the ordered differences \(\alpha_k-\alpha_j\)
\((j,k\in[n],\ j\neq k)\) are pairwise distinct.
\end{lemma}

\begin{proof}
For \(n=1\), take \(\alpha_1=0\).  Otherwise we can always choose an odd prime
\(p\in[n,2n)\) by Bertrand's
postulate.  Let \(r_j\) be the least residue of \((j-1)^2\) modulo
\(p\), and set
\[
 \alpha_j:=2p(j-1)+r_j.
\]
This is the classical Erd{\H{o}}s--Tur{\'a}n construction,
which ensures that all nonzero ordered differences \(\alpha_k-\alpha_j\) are
pairwise distinct~\cite{ErdosTuran41}.
The construction runs in polynomial time.  Finally,
\(p<2n\) and \(r_n<p\) give \(\alpha_n=O(n^2)\), while \(p\geq n\) gives
\[
\alpha_n\geq2p(n-1)\geq2n(n-1).
\]
Hence \(\alpha_n=\Theta(n^2)\) for \(n\geq2\).
\end{proof}

Set
\[
 c:=\alpha_n, \quad L:=2c+1,\quad \rho_i:=c+(i-1)L\quad(i\in[M]).
\]
We call \(\rho_i\) the \emph{row coordinate} of row \(i\).  For
\(A_{ij}=1\) and \(k\neq j\), the position
\(\rho_i-\alpha_j+\alpha_k\) is a
\emph{cross-term coordinate}.  Set
\[
 W:=\sum_{i\in[M]}\sum_{j\in[n]}A_{ij},
 \qquad
 \nu:=(n-1)W.
\]

Let \(N\) be the least power of two greater than \(\rho_M+c\), and
work in \(R_N^-:=\Z[x]/(x^N+1)\).  We use degree-\(<N\) representatives
and their coefficient norms.  Define
\begin{align*}
 h(x)&:=\sum_{A_{ij}=1}x^{\rho_i-\alpha_j},\\
 t(x)&:=2\sum_{i=1}^M b_i x^{\rho_i}
       +\sum_{\substack{A_{ij}=1\\k\neq j}}
          x^{\rho_i-\alpha_j+\alpha_k},\\
 a_{\boldsymbol\xi}(x)&:=\sum_{k=1}^n\xi_kx^{\alpha_k}.
\end{align*}

\begin{lemma}
\label{lem:boolean-identity}
For every Boolean vector
\(\boldsymbol\xi\in\{0,1\}^n\),
\[
 \|2h(x)a_{\boldsymbol\xi}(x)
       -t(x)\|_2^2
 =\nu+4\|\mathbf A\boldsymbol\xi-\mathbf b\|_2^2.
\]
\end{lemma}

\begin{proof}
Note that
\[
2h(x)a_{\boldsymbol{\xi}}(x)
= 2\sum_{A_{ij}=1}\sum_{k=1}^n \xi_k x^{\rho_i - \alpha_j + \alpha_k}.
\]
Splitting the inner sum according to $k=j$ and $k\neq j$ yields
\begin{equation}
2h(x)a_{\boldsymbol{\xi}}(x)
= 2\sum_{A_{ij}=1} \xi_j x^{\rho_i}
+ 2\sum_{\substack{A_{ij}=1\\k\neq j}} \xi_k x^{\rho_i - \alpha_j + \alpha_k}.
\label{equ::ha}
\end{equation}

Subtracting $t(x)$ from \eqref{equ::ha} gives
\begin{equation}
2h(x)a_{\boldsymbol{\xi}}(x) - t(x)
= 2\sum_{i=1}^M\Bigl(\sum_{j:A_{ij}=1}\xi_j - b_i\Bigr)x^{\rho_i}
+ \sum_{\substack{A_{ij}=1\\k\neq j}}(2\xi_k - 1)x^{\rho_i - \alpha_j + \alpha_k}.
\label{equ::sol}
\end{equation}
Because $(\mathbf{A}\boldsymbol{\xi})_i=\sum_{j:A_{ij}=1}\xi_j$, the coefficient of $x^{\rho_i}$ in~(\ref{equ::sol}) is precisely $2((\mathbf{A}\boldsymbol{\xi})_i-b_i)$.

To compute the \(\ell_2\)-norm of~\eqref{equ::sol}, we must verify that all its monomials have distinct exponents.
The row coordinates \(\rho_i\) are pairwise distinct because \(L>0\).

\textit{(i) Cross-term coordinates do not equal row
coordinates.}
A row coordinate is some
$\rho_{i'}=c+(i'-1)L$. A cross-term coordinate is
$\rho_i-\alpha_j+\alpha_k$ with $A_{ij}=1$ and $k\neq j$. Equality of the two
would require $\rho_{i'}-\rho_i=\alpha_k-\alpha_j$. If $i'=i$ this forces
$\alpha_j=\alpha_k$, hence $j=k$, contradicting $k\neq j$. If $i'\neq i$
then $|\rho_{i'}-\rho_i|\ge L=2c+1$, whereas $|\alpha_k-\alpha_j|\le c$
because $0\le\alpha_j,\alpha_k\le c$.
Thus equality is impossible.

\textit{(ii) Cross-term coordinates for a fixed row are
distinct.}
Fix $i$ and suppose
\[
\rho_i-\alpha_j+\alpha_k = \rho_i-\alpha_{j'}+\alpha_{k'}
\quad\text{with }k\neq j,\; k'\neq j'.
\]
Cancelling $\rho_i$ gives $\alpha_k-\alpha_j=\alpha_{k'}-\alpha_{j'}$. By
Lemma~\ref{lem:sidon} all nonzero ordered differences $\alpha_k-\alpha_j$ are distinct,
so $(j,k)=(j',k')$.

\textit{(iii) Cross-term coordinates for different rows are distinct.}
All cross-term coordinates for row \(i\) lie in
\([\rho_i-c,\rho_i+c]=[(i-1)L,iL-1]\).  These intervals are disjoint for
distinct rows.

Since $N$ is a power of two larger than $\rho_M+c$, all exponents lie in
$[0,N)$ and no wrap-around modulo $x^N+1$ occurs.  Consequently, the two sums
in~\eqref{equ::sol} have disjoint supports.  The second sum contains exactly
$\nu=(n-1)W$ monomials, one for each triple $(i,j,k)$ with $A_{ij}=1$ and
$k\neq j$.  Because $\xi_k\in\{0,1\}$, each coefficient satisfies
$(2\xi_k-1)^2=1$.  Therefore
\[
\|2h(x)a_{\boldsymbol{\xi}}(x)-t(x)\|_2^2
=4\|\mathbf{A}\boldsymbol{\xi}-\mathbf{b}\|_2^2+\nu.
\]
\end{proof}

\subsection{Controlling every ideal element}

Lemma~\ref{lem:boolean-identity} applies only to the
polynomials \(a_{\boldsymbol\xi}\) encoding Boolean vectors.  We now add
coefficients at odd powers of \(y\) so that every \(q\in R_d\) not of the
form \(a_{\boldsymbol\xi}(y^2)\) with
\(\boldsymbol\xi\leq\boldsymbol\mu\) satisfies
\(\|gq-T_{\boldsymbol\mu}\|_2^2>\Delta_{\boldsymbol\mu}+4\).

\begin{lemma}
The polynomial \(h(x)=\sum_{A_{ij}=1}x^{\rho_i-\alpha_j}\) has exactly
\(W\) distinct monomials.  Equivalently, the exponents
\(\rho_i-\alpha_j\) with \(A_{ij}=1\) are pairwise distinct.
\end{lemma}

\begin{proof}
Suppose \((i,j)\neq(i',j')\), \(A_{ij}=A_{i'j'}=1\), and
\[
\rho_i-\alpha_j=\rho_{i'}-\alpha_{j'}.
\]

\textit{Case 1: \(i=i'\).}
Then \(\alpha_j=\alpha_{j'}\), so strict monotonicity of the
\(\alpha_j\)'s gives \(j=j'\), a contradiction.

\textit{Case 2: \(i\neq i'\).}
The spacing of the row coordinates gives
\[
|\rho_{i'}-\rho_i|\geq L=2c+1,
\qquad
|\alpha_{j'}-\alpha_j|\leq c,
\]
which contradicts
\(\rho_{i'}-\rho_i=\alpha_{j'}-\alpha_j\).

Finally, \(0\leq\rho_i-\alpha_j\leq\rho_M<N\), so no wrap-around occurs
in \(R_N^-\).  Hence all \(W\) exponents are distinct.
\end{proof}





Since these \(W\) monomials have unit coefficients, \(\|h\|_1=W\).  Set
\[
 s:=\lceil\sqrt{n+4}\rceil,
 \qquad
 P:=1+s\bigl(sW+\lceil\|t\|_2\rceil
                    +\lceil\sqrt{\nu+4}\rceil\bigr).
\]
Let \(d:=2N\) and identify \(R_N^-\) with the even-coordinate subring of
\(R_d=\Z[y]/(y^d+1)\) via \(x\mapsto y^2\).  For a mask
\(\boldsymbol\mu\in\{0,1\}^n\), put
\[
 u_{\boldsymbol\mu}(x):=\sum_{k=1}^n\mu_kx^{\alpha_k}
\]
and define
\begin{align}
 g(y)&:=2h(y^2)+2Py,                                      \label{eq:generator}\\
 T_{\boldsymbol\mu}(y)&:=t(y^2)+Py\,u_{\boldsymbol\mu}(y^2),
                                                             \label{eq:masked-target}\\
 \Delta_{\boldsymbol\mu}&:=P^2\|\boldsymbol\mu\|_0+\nu.
\label{eq:masked-radius}
\end{align}

Consider the principal ideal generated by $g$.  Every \(q\in R_d\) has a unique form
\(q=a(y^2)+yb(y^2)\) with \(a,b\in R_N^-\).  Direct multiplication gives
\begin{equation}
\label{eq:block-product}
 gq-T_{\boldsymbol\mu}
 =(2ha+2Pxb-t)(y^2)+y(2hb+2Pa-Pu_{\boldsymbol\mu})(y^2),
\end{equation}
where both \(2ha+2Pxb-t\) and
\(2hb+2Pa-Pu_{\boldsymbol\mu}\) belong to \(R_N^-\).


\begin{lemma}
\label{lem:admissible}
If \(\boldsymbol\xi\in\{0,1\}^n\) and
\(\boldsymbol\xi\leq\boldsymbol\mu\), then
\begin{equation}
\label{eq:allowed-distance}
 \|g(y)a_{\boldsymbol\xi}(y^2)
       -T_{\boldsymbol\mu}(y)\|_2^2
 =\Delta_{\boldsymbol\mu}
  +4\|\mathbf A\boldsymbol\xi-\mathbf b\|_2^2.
\end{equation}
Every $q\in R_d$ not of the form \(a_{\boldsymbol\xi}(y^2)\) for some
\(\boldsymbol\xi\in\{0,1\}^n\) with
\(\boldsymbol\xi\leq\boldsymbol\mu\) satisfies
\[
 \|g(y)q(y)
       -T_{\boldsymbol\mu}(y)\|_2^2
 >\Delta_{\boldsymbol\mu}+4.
\]
\end{lemma}

\begin{proof}
Let \(\omega:=\|\boldsymbol\mu\|_0\) be the Hamming weight of
\(\boldsymbol\mu\).

\textit{Case 1: $q=a_{\boldsymbol\xi}(y^2)$ for a Boolean vector
\(\boldsymbol\xi\leq\boldsymbol\mu\).}
Equation~\eqref{eq:block-product} gives
\begin{equation}
\label{eq:case-one-blocks}
 g(y)q(y)-T_{\boldsymbol\mu}(y)
 =\underbrace{2h(y^2)a_{\boldsymbol\xi}(y^2)-t(y^2)}_{\text{even block}}
 +
 \underbrace{Py\bigl(2a_{\boldsymbol\xi}(y^2)
                   -u_{\boldsymbol\mu}(y^2)\bigr)}_{\text{odd block}}.
\end{equation}

\paragraph{The even block.}
The even block in~\eqref{eq:case-one-blocks} is supported on
the even coordinates, and the substitution \(x=y^2\) preserves its
coefficient \(\ell_2\)-norm. Hence
Lemma~\ref{lem:boolean-identity} gives
\[
 \|2h(y^2)a_{\boldsymbol\xi}(y^2)-t(y^2)\|_2^2
 =\|2h(x)a_{\boldsymbol\xi}(x)-t(x)\|_2^2
 =\nu+4\|\mathbf A\boldsymbol\xi-\mathbf b\|_2^2.
\]

\paragraph{The odd block.}
The odd block in~\eqref{eq:case-one-blocks} is
\[
 Py\bigl(2a_{\boldsymbol\xi}(y^2)-u_{\boldsymbol\mu}(y^2)\bigr)
 =P\sum_{k=1}^n(2\xi_k-\mu_k)y^{2\alpha_k+1}.
\]
For every $k$, the conditions
\(\boldsymbol\xi\leq\boldsymbol\mu\) and
\(\boldsymbol\xi,\boldsymbol\mu\in\{0,1\}^n\) imply
\((2\xi_k-\mu_k)^2=\mu_k\).  Moreover,
\(2\alpha_k+1\leq2c+1\leq N<d\), so the corresponding odd-coordinate
positions are distinct and do not wrap around in $R_d$.
Hence the odd block in~\eqref{eq:case-one-blocks} has squared norm
\[
 P^2\sum_{k=1}^n(2\xi_k-\mu_k)^2
 =P^2\|\boldsymbol\mu\|_0=P^2\omega.
\]
Since the even and odd coordinates are disjoint, adding the two contributions
proves~\eqref{eq:allowed-distance}.

\textit{Case 2: all remaining $q\in R_d$.}
Write $q(y)=a(y^2)+yb(y^2)$ with $a,b\in R_N^-$, and set
\[
 \Phi(a,b):=\|b\|_2^2+
 \left\|a-\frac{u_{\boldsymbol\mu}}2\right\|_2^2.
\]
For \(0\leq r<N\), write \(a_r\) and \(b_r\) for the
coefficients of \(x^r\) in \(a\) and \(b\), respectively.
For each $k$ with \(\mu_k=1\), integrality gives
\((a_{\alpha_k}-1/2)^2\geq1/4\), with equality exactly when
\(a_{\alpha_k}\in\{0,1\}\).
Each coefficient \(a_r\) with
\(r\notin\{\alpha_k:\mu_k=1\}\) contributes \(a_r^2\) to
\(\Phi(a,b)\), and each coefficient \(b_r\) contributes \(b_r^2\).
Consequently,
\[
 \Phi(a,b)=\frac\omega4
\]
exactly when $b=0$ and
\(a(x)=a_{\boldsymbol\xi}(x)\) for some Boolean vector
\(\boldsymbol\xi\leq\boldsymbol\mu\).  For every other pair,
\(\Phi(a,b)\) is larger by at least \(1\).  Hence
\[
 \Phi(a,b)\geq\frac{\omega+4}{4}.
\]

Identify a polynomial in \(R_d\) with the ordered pair of its
even- and odd-degree coefficient blocks, equipped with the product Euclidean
norm.
Equation~\eqref{eq:block-product} becomes
\[
 gq-T_{\boldsymbol\mu}
 =2P\bigl((xb,a)-(0,u_{\boldsymbol\mu}/2)\bigr)
  +2(ha,hb)-(t,0).
\]
Multiplication by $x$ is an isometry in $R_N^-$, so
\(\|xb\|_2=\|b\|_2\).  Young's inequality and
\(\|h\|_1=W\) give the first bound below, while the triangle inequality gives
the second:
\[
 \|(ha,hb)\|_2\leq W\|(a,b)\|_2,
 \qquad
 \|(a,b)\|_2\leq\sqrt{\Phi(a,b)}+\frac{\sqrt\omega}{2}.
\]
The reverse triangle inequality therefore yields
\[
 \|gq-T_{\boldsymbol\mu}\|_2
 \geq2(P-W)\sqrt{\Phi(a,b)}-W\sqrt\omega-\|t\|_2.
\]
Since $P\geq1+s^2W>W$, the lower bound on \(\Phi(a,b)\) gives
\[
 \|gq-T_{\boldsymbol\mu}\|_2
 \geq P\sqrt{\omega+4}
   -W\bigl(\sqrt{\omega+4}+\sqrt\omega\bigr)-\|t\|_2.
\]

Because $0\leq\omega\leq n$ and
\(s=\lceil\sqrt{n+4}\rceil\),
\[
 \sqrt{\omega+4}-\sqrt\omega\geq\frac2s,
 \qquad
 \sqrt{\omega+4}+\sqrt\omega\leq2s.
\]
The definition of $P$ also gives the strict inequality
\[
 \frac{2P}{s}
 =\frac2s+2sW+2\lceil\|t\|_2\rceil
   +2\lceil\sqrt{\nu+4}\rceil
 >2sW+\|t\|_2+\sqrt{\nu+4}.
\]
Combining these estimates gives
\[
\begin{aligned}
 \|gq-T_{\boldsymbol\mu}\|_2-P\sqrt\omega
 &\geq P\bigl(\sqrt{\omega+4}-\sqrt\omega\bigr)
   -W\bigl(\sqrt{\omega+4}+\sqrt\omega\bigr)-\|t\|_2 \\
 &>\sqrt{\nu+4}.
\end{aligned}
\]
Squaring and discarding the nonnegative cross term now gives
\[
 \|gq-T_{\boldsymbol\mu}\|_2^2
 >\bigl(P\sqrt\omega+\sqrt{\nu+4}\bigr)^2
 \geq P^2\omega+\nu+4
 =\Delta_{\boldsymbol\mu}+4.
\]
This proves the second assertion.
\end{proof}

The preceding lemma yields the reduction for masked Boolean
systems.  Its X3C specialization gives the main hardness result, while the
fact that the ideal is independent of the mask will be used for the
fixed-family construction in Section~\ref{sec:fixed}.
For a full-rank principal ideal \(I=(g)\lhd R_d\), we use
\((I,T,\Delta)\) as shorthand for the standard CVP input
\((\mathbf B_g,\coeff(T),\Delta)\).

\begin{theorem}
\label{thm:compiler}
There is a polynomial-time reduction from the masked Boolean system
\eqref{eq:masked-system} to exact decision-CVP that maps
\((\mathbf A,\mathbf b,\boldsymbol\mu)\) to
\((I,T_{\boldsymbol\mu},\Delta_{\boldsymbol\mu})\), where
\(g\in R_d\setminus\{0\}\) and
\(I\mathrel{:=}(g)\lhd R_d\) is a full-rank principal ideal with a multiplication
basis depending only on \((\mathbf A,\mathbf b)\).
Here \(d=O((M+1)n^2)\) is a power of two, and
\(T_{\boldsymbol\mu}\) and \(\Delta_{\boldsymbol\mu}\) are integral.
\end{theorem}

\begin{proof}
Equation~\eqref{eq:allowed-distance} shows that the element
\(g a_{\boldsymbol\xi}(y^2)\) associated with a Boolean vector
\(\boldsymbol\xi\leq\boldsymbol\mu\) has squared distance
\(\Delta_{\boldsymbol\mu}\) from \(T_{\boldsymbol\mu}\)
precisely when
\(\mathbf A\boldsymbol\xi=\mathbf b\), i.e., when \(\boldsymbol\xi\) is a
feasible solution of the masked Boolean system~\eqref{eq:masked-system}.
If \(\mathbf A\boldsymbol\xi\neq\mathbf b\), then the integrality of
\(\mathbf A\boldsymbol\xi-\mathbf b\) implies
\(\|\mathbf A\boldsymbol\xi-\mathbf b\|_2^2\geq1\), so the squared distance
is at least \(\Delta_{\boldsymbol\mu}+4\).
The second assertion
of Lemma~\ref{lem:admissible} gives the same lower bound for
every remaining \(q\in R_d\).  Consequently,
\[
\begin{aligned}
 {\dist(T_{\boldsymbol\mu},I)^2}
 &=\Delta_{\boldsymbol\mu} &&\text{in a YES instance},\\
 {\dist(T_{\boldsymbol\mu},I)^2}
 &\geq\Delta_{\boldsymbol\mu}+4 &&\text{in a NO instance}.
\end{aligned}
\]

The coefficient of \(y\) in \(g\) is \(2P\neq0\), so \(g\neq0\).
Since \(R_d\) is an integral domain, multiplication by \(g\) is injective.
Hence \((g)\) has full \(\mathbb Z\)-rank.  Moreover, the elements \(v\in(g)\)
satisfying \(\|v-T_{\boldsymbol\mu}\|_2^2\leq\Delta_{\boldsymbol\mu}\) are in
bijection with the solutions of \eqref{eq:masked-system}.

Finally, \(c=O(n^2)\) and
\(\rho_M+c=O((M+1)n^2)\), so \(d=O((M+1)n^2)\).  We also have
\(W\leq Mn\), \(\nu\leq Mn^2\), and
\(\|t\|_2^2=4\|\mathbf b\|_0+\nu\).  Hence every constructed quantity has
polynomial encoding length, and the reduction runs in polynomial time.
The generator \(g\in R_d\setminus\{0\}\) and its
multiplication basis depend only on
\((\mathbf A,\mathbf b)\), whereas \(T_{\boldsymbol\mu}\) and
\(\Delta_{\boldsymbol\mu}\) also depend on the mask \(\boldsymbol\mu\).
\end{proof}

We now apply Theorem~\ref{thm:compiler} with
\((\mathbf A,\mathbf b,\boldsymbol\mu)
  \mathrel{:=}(\mathbf H,\mathbf 1_m,\mathbf 1_n)\).

\begin{theorem}
\label{thm:cyclotomic-cvp}
Exact decision-CVP is \(\NP\)-complete on the coefficient lattices of
full-rank principal ideals in the power-of-two cyclotomic rings \(R_d\).
The hard instances have squared-distance gap \(\Delta\) versus
\(\Delta+4\).
\end{theorem}

\begin{proof}
Theorem~\ref{thm:compiler} gives a deterministic reduction from X3C and the
claimed gap.  For membership in \(\NP\), a YES certificate is
\(\coeff(v)\) for an element \(v\in(g)\) satisfying
\(\|v-T\|_2^2\leq\Delta\).  Write
\(\coeff(v)=(v_0,\ldots,v_{d-1})^\top\) and
\(\coeff(T)=(T_0,\ldots,T_{d-1})^\top\).  Each coefficient satisfies
\(|v_r|\leq|T_r|+\lceil\sqrt\Delta\rceil\), and membership in \((g)\) is
verified by solving the nonsingular multiplication system
\(\mathbf B_g\mathbf z=\coeff(v)\) and checking
\(\mathbf z\in\Z^d\).
The verifier also checks \(\|v-T\|_2^2\leq\Delta\).

\end{proof}

The same characterization of the elements satisfying the
distance threshold shows that exact search-CVP on these ideals is
\(\NP\)-hard under polynomial-time Turing reductions.
Indeed, one call to an exact search-CVP oracle returns a closest lattice vector
\(\coeff(v)\) for some \(v\in(g)\).
If \(\|v-T\|_2^2\leq\Delta\), solving
\(\mathbf B_g\mathbf z=\coeff(v)\) recovers \(q\) through
\(\mathbf z=\coeff(q)\).  The coefficients of \(q\) at
\(y^{2\alpha_k}\) give the exact cover.  Otherwise the X3C instance is
negative.

By Lemma~\ref{lem:canonical}, both results hold under the
canonical embedding.  Squared thresholds and gaps are multiplied by \(d\),
and the closest elements are unchanged.

\subsection{Cyclic extension}
\label{sec:cyclic-cvp}

Let \(S_D:=\Z[X]/(X^D-1)\), identified with \(\Z^D\) by coefficient vectors.
Multiplication by \(X\) rotates the coordinates, so every full-rank ideal of
\(S_D\) is a \emph{cyclic lattice}.

The following lemma lifts the principal cyclotomic instances to
principal cyclic ideals in \(S_{2d}\).
Here \(I\) is the
original principal ideal, \(T\) is its target, and \(H\) is the
squared-distance threshold. Under the lift, squared distances
are multiplied by \(8\), and points of \(I\) within squared distance \(H\)
of \(T\) correspond bijectively to points of the lifted ideal within
squared distance \(8H\) of the lifted target.

\begin{lemma}
\label{lem:principal-crt}
Let \(d\) be a power of two, \(I=(g)\lhd R_d\) with \(g\neq0\), and
\(T\in R_d\).  Fix \(H\geq0\), choose an integer \(B\) with \(B^2>4H\),
and set \(D:=2d\).  Let \(\iota_d:R_d\to S_D\) copy the
degree-\(<d\) representative, and define
\[
 G:=(1-X^d)\iota_d(g)+B(1+X^d),
 \qquad
 \widehat T:=2(1-X^d)\iota_d(T).
\]
Then \(J:=(G)\) is full rank, and
\[
 \|2(1-X^d)\iota_d(v)-\widehat T\|_2^2
 =8\|v-T\|_2^2
 \qquad(v\in I).
\]
Every \(w\in J\) satisfying
\(\|w-\widehat T\|_2^2\leq8H\) arises uniquely in this way.
\end{lemma}

\begin{proof}
For an arbitrary \(Q\in S_D\), write \(Q=p+X^dq\), where
\(p,q\in\mathbb Z[X]\) have degree less than \(d\).  Put \(r:=p+q\) and
\(z:=p-q\).  We view these degree-\(<d\)
polynomials as elements of \(R_d\) in products with \(g\).  Equivalently,
\begin{equation}
\label{eq:q-rz-decomposition}
 Q=\frac12\bigl((1+X^d)r+(1-X^d)z\bigr).
\end{equation}
If
\(F:=GQ-\widehat T=F_0+X^dF_1\), direct expansion in \(R_d\) gives
\begin{equation}
\label{eq:principal-crt-distance}
\begin{gathered}
 F_0=Br+gz-2T,\qquad F_1=Br-(gz-2T),\\
 \|F\|_2^2=\|F_0\|_2^2+\|F_1\|_2^2
 =2\|gz-2T\|_2^2+2B^2\|r\|_2^2.
\end{gathered}
\end{equation}
If \(r\neq0\), then the integer coefficient vector of \(r\)
has squared norm at
least one, so \eqref{eq:principal-crt-distance} is at least \(2B^2>8H\).
Hence \(\|GQ-\widehat T\|_2^2\leq8H\) implies
\(r=0\). Then \(q=-p\) and \(z=2p\). Substituting into
\eqref{eq:q-rz-decomposition} gives
\[
 Q=\frac12(1-X^d)z=(1-X^d)p.
\]
Using \((1-X^d)(1+X^d)=0\) and
\((1-X^d)^2=2(1-X^d)\), we first get
\(GQ=2(1-X^d)\iota_d(g)p\).  By the definition of \(\iota_d(gp)\), the
difference \(\iota_d(g)p-\iota_d(gp)\) is divisible by \(X^d+1\).  Since
\((1-X^d)(X^d+1)=1-X^{2d}=0\) in \(S_D\), we obtain
\[
 GQ=2(1-X^d)\iota_d(gp),
 \qquad
\|GQ-\widehat T\|_2^2=8\|gp-T\|_2^2.
\]
Thus \(v=gp\) maps to \(2(1-X^d)\iota_d(v)\).  The coefficients
of degrees \(0,\ldots,d-1\) in this image equal those of \(2v\), so \(v\) is
uniquely determined.

Over \(\Q\), the Chinese remainder theorem gives
\[
 S_D\otimes_{\Z}\Q
 \cong \Q[X]/(X^d-1)\times\Q[X]/(X^d+1),
 \qquad
 G\longmapsto(2B,2g).
\]
Since \(d\) is a power of two, \(X^d+1\) is irreducible over \(\Q\), so
\(\Q[X]/(X^d+1)\) is a field.  Both components are units: \(2B\) is a
nonzero scalar in the first factor, and \(2g\neq0\) in the second.  Hence
multiplication by \(G\) is invertible over \(\Q\), and \(J\) has full rank.
\end{proof}

\begin{theorem}
\label{thm:cyclic-cvp}
Exact decision-CVP for tuples
\((D,G,\widehat T,\widehat\Delta)\), where \(D\) is a power of two,
{\(G,\widehat T\in S_D\)}, \((G)\lhd S_D\) is full rank, and
{\(\widehat\Delta\in\Z_{\geq0}\)}, is
\(\NP\)-complete.  The lattice basis is computed from the
\(D\) cyclic shifts of \(G\).  The reduction has squared-distance gap
\(\widehat\Delta\) versus \(\widehat\Delta+32\).
\end{theorem}

\begin{proof}
Apply Lemma~\ref{lem:principal-crt} to the instance from
Theorem~\ref{thm:cyclotomic-cvp} with
\[
 H:=\Delta+4,
 \qquad
 B:=2\lceil\sqrt H\rceil+1,
 \qquad
 \widehat\Delta:=8\Delta.
\]
The closest squared distance equals \(\widehat\Delta\) in a YES instance and
is at least \(\widehat\Delta+32\) in a NO instance.  All lifted data have
polynomial encoding length.
Membership in \(\NP\) follows as in
Theorem~\ref{thm:cyclotomic-cvp}, using the cyclic shifts of \(G\) as a
basis.
\end{proof}

\begin{remark}
Theorem~\ref{thm:cyclic-cvp} settles one of Micciancio's open questions:
exact Euclidean decision-CVP is \(\NP\)-hard on cyclic
lattices~\cite{Micciancio2007}.
The lift also gives \(\NP\)-hardness of
exact search-CVP on the same cyclic ideals under polynomial-time Turing
reductions.  Given a closest \(\widehat v\in J\), if
\(\|\widehat v-\widehat T\|_2^2>\widehat\Delta\), the X3C instance is
negative.  Suppose instead that
\(\|\widehat v-\widehat T\|_2^2\leq\widehat\Delta\).
Lemma~\ref{lem:principal-crt} then gives a unique \(v\in I\) with
\(\widehat v=2(1-X^d)\iota_d(v)\). The coefficients of \(\widehat v\) in
degrees \(0,\ldots,d-1\) are twice the corresponding coefficients of \(v\).
Thus they determine \(\coeff(v)\). Solving
\(\mathbf B_g\mathbf z=\coeff(v)\) now yields
\(\mathbf z=\coeff(q)\). The coefficients of \(q\) at
\(y^{2\alpha_k}\) give the exact cover.
\end{remark}
